\subsection{Ponderación disponible y orden de prioridad}

La técnica de votación ponderada prevista en el kit de formatos reparte 100 puntos entre veinte requisitos \cite{wiegers2013}. Sin embargo, la información entregada para esta entrega consiste en \textbf{una sola tabla de ponderación} con valores enteros entre 3 y 5, no en tres boletas de 100 puntos. Por ello se registró esa tabla tal cual, sin generar votantes adicionales, y se calcula la correlación con los valores disponibles. Esta limitación se declara explícitamente en la Sección~\ref{sec:comparacion}.

\begin{table}[ht]
\centering
\footnotesize
\caption{Ponderación real entregada (una tabla, escala 3--5).}
\label{tab:votacion_real}
\resizebox{\linewidth}{!}{%
\begin{tabular}{@{}llr@{}}
\toprule
ID & Requisito & Puntos \\
\midrule
R01 & Reservar citas médicas. & 5 \\
R02 & Enviar recordatorios por mensaje de texto. & 4 \\
R03 & La aplicación debe ser rápida. & 5 \\
R04 & Registrar la llegada de pacientes. & 4 \\
R05 & Desarrollar la plataforma en Java 21 y PostgreSQL. & 4 \\
R06 & Consultar pacientes del día y marcar consultas atendidas. & 5 \\
R07 & Proteger los datos de salud de los pacientes. & 5 \\
R08 & Funcionar en teléfonos Android 10 o superior. & 3 \\
R09 & Registrar cancelaciones tardías como inasistencias. & 4 \\
R10 & Ser fácil de utilizar para adultos mayores. & 5 \\
R11 & Entregar el proyecto en seis meses. & 3 \\
R12 & Cerrar automáticamente la sesión después de 10 minutos. & 4 \\
R13 & Registrar llegada e imprimir tickets en el quiosco. & 4 \\
R14 & Mantener la plataforma siempre disponible. & 5 \\
R15 & Generar reportes mensuales de inasistencias. & 4 \\
R16 & Validar cada incremento cada dos semanas. & 3 \\
R17 & Cumplir la Ley de Protección de Datos Personales. & 5 \\
R18 & Integrarse con el sistema de facturación actual. & 5 \\
R19 & Recuperar el servicio en menos de 30 minutos sin perder citas. & 5 \\
R20 & Revisar los cambios de código antes de integrarlos. & 3 \\
\bottomrule
\end{tabular}%
}
\end{table}

El total de la tabla es 85 puntos. Los requisitos con 5 puntos son R01, R03, R06, R07, R10, R14, R17, R18 y R19; los de 4 puntos son R02, R04, R05, R09, R12, R13 y R15; y los de 3 puntos son R08, R11, R16 y R20. Los datos se registran en \texttt{datos/votacion.csv}. El procesamiento se realiza con \texttt{datos/calcular\_spearman.py}.
